% Section III, organised as the four stages of Fig. 1 (fig1_architecture_v2.pdf).
% Drop-in for the Overleaf draft: uses its macros (\tool, \dslkw, \dsllit, \hl)
% and its figure label. The \providecommand lines are no-ops where the macros
% already exist. Every subsection = one stage of the figure, in the figure's
% order: what comes in, the components (named as the boxes), what goes out.
\providecommand{\tool}{REAL}
\providecommand{\dslkw}[1]{\texttt{#1}}
\providecommand{\dsllit}[1]{\texttt{#1}}
\providecommand{\hl}[1]{\textcolor{red}{#1}}
\providecommand{\todo}[1]{\textcolor{red}{[TODO: #1]}}

\section{The \tool{} Tool}
\label{sec:tool}

\begin{figure*}[t]
  \centering
  \includegraphics[width=\textwidth]{IMG/fig1_architecture.pdf}
  \caption{Overview of \tool{}. (1)~The engineer
  specifies $R_0$, which \tool{} parses and checks for well-formedness and for
  assumptions about the system under test. (2)~\tool{}
  derives test scenarios from the requirement's grammar by grammatical evolution
  or exhaustive enumeration and executes them in CARLA via Scenic, with the
  learned detector and the braking controller in the loop.
  (3)~\tool{} classifies each failure as valid or
  spurious against the domain assumptions of $R_0$ and abstracts valid failures
  into obstacles with supporting evidence. (4)~The
  engineer triages obstacles and assumptions, and \tool{} derives $R_1$ with
  each change attributed to the system~[S], the domain~[D] or the
  requirement~[R]. System changes trigger a new round; domain and requirement
  changes are re-diagnosed on the recorded traces, without re-simulation.
  Robots mark the activities that \tool{} automates and that the REAL framework
  performs manually; hard hats mark the engineer's decisions. Starred components are contributions of this paper; the others are reused from REAL~\cite{real2026} or from third-party tools~\cite{carla2017,scenic2023,grape2022}.}
  \label{fig:real-architecture}
\end{figure*}

\tool{} automates the REAL identify--analyse--mitigate
loop~\cite{real2026}: each iteration refines a requirement $R_0$ into $R_1$ in
four stages (Fig.~\ref{fig:real-architecture}), and $R_1$ seeds the next
iteration. The engineer intervenes at two points only: specifying $R_0$
(stage~1) and interpreting the diagnosed failures (stage~4). All intermediate
stages are automated, and every computed result is presented together with its
supporting evidence. Consecutive stages communicate through a single artefact
each: the checked requirement, the execution traces, the diagnosis, and the
signed decisions. Section~\ref{sec:features} details the features, F1--F14,
that realise each stage.

\subsection{Stage 1: Requirement Elicitation}
\label{sec:elicit}
The engineer specifies $R_0$ in the REAL DSL~\cite{real2026} through a
structured editor (F1). The REAL-DSL parser validates each edit and renders the
goal model; the requirement check rejects specifications that \tool{} cannot
execute or whose assumptions constrain the system under test (F1, F2). The stage
outputs the checked requirement $R_0$.

\subsection{Stage 2: Generate and Execute Tests}
\label{sec:generate}
\tool{} derives test scenarios from the requirement's grammar, by exhaustive
enumeration or by grammatical evolution~\cite{oneill2001,grape2022}, and
executes them in CARLA~\cite{carla2017} via Scenic~\cite{scenic2023}, with the
system under test in the loop. A safety monitor scores every run and a recorder
writes its per-step trace (F3, F4). The stage runs as a job on a GPU cluster
and outputs the execution traces.

\subsection{Stage 3: Diagnose Test Results}
\label{sec:diagnose}
\tool{} judges each run valid or spurious against the assumptions of $R_0$
(F2), types each valid failure from its trace and abstracts recurring failures
into obstacles with supporting evidence (F5), and clusters the traces by their
event sequences to expose behaviour patterns that the taxonomy does not name
(F6). The stage outputs the diagnosis:
valid failures, obstacles, and a verdict on each assumption.

\subsection{Stage 4: Decide and Refine}
\label{sec:decide}
In the review interface, the engineer triages each obstacle and assumption,
supported by the tool's suggestions and evidence (F7). \tool{} then derives
$R_1$, attributing each change to the system~[S], the domain~[D] or the
requirement~[R], and records the signed decisions (F8). System changes trigger
a new iteration from stage~1; domain and requirement changes are re-diagnosed on
the recorded traces, without re-simulation (F8).

\subsection{Implementation}
\label{sec:impl}
\tool{} comprises a Python back end and a React/TypeScript front end. The back
end provides the REAL-DSL parser, a FastAPI service that orchestrates Scenic and
CARLA, the analysis pipeline, and job control on a Slurm cluster through an
Apptainer image. Every run is persisted as a self-describing folder (metadata with
the hashes of the scene template and grammar and the scene settings,
per-simulation tables, and per-step traces), so that every analysis can be reproduced from the recorded traces
without re-simulation. A Docker image packages the review interface and the requirement check with
recorded runs, and requires neither a GPU nor a simulator. The implementation is
covered by 585 automated tests (317 for the back end and 268 for the
front end), none of which requires the simulator.

% ---------------------------------------------------------------------------
\section{Key Features}
\label{sec:features}

REAL~\cite{real2026} establishes the identify--analyse--mitigate loop and
demonstrates it on an autonomous braking system, with the central analysis
steps performed by hand: \emph{``failure grouping and obstacle identification
were performed primarily through manual inspection of traces''}, which the
authors identify as a natural candidate for automation~\cite[Sec.~V-D]{real2026}.
\tool{} realises this automation through the features F1--F14, each
implemented by the components marked $\bigstar$ in
Fig.~\ref{fig:real-architecture}.

% F1 <- inventory F5-F7, F13-F21; F2 <- F52-F56, F59-F60; F3 <- F30-F39, F44-F49;
% F4 <- F61-F66; F5 <- F67-F81, F84; F6 <- F82-F83, F98; F7 <- F89-F119; F8 <- F122-F139.
\paragraph*{F1 Entailment-aware specification (stage 1)} \tool{} extends the
REAL DSL with an \dslkw{ensuring} slot for soft goals, so a requirement states
the system (\dslkw{performed by}), the domain (\dslkw{assuming}) and the
requirement (\dslkw{ensuring}) separately. The grammar-aware editor completes
only what the executor implements (modules, measurable quantities with units,
operators), links each phrase to its node in the live KAOS goal model, and
refuses unknown module names instead of silently substituting another system.

\paragraph*{F2 Admissibility derived from the requirement (stages 1, 3)}
Where REAL states the admissibility predicate $\Phi_{\mathit{valid}}$ apart
from the requirement, \tool{} derives it from the \dslkw{assuming} clauses.
Each clause of the form \emph{quantity operator value} is matched against a
vocabulary of measurable world quantities and checked, for every run, against
either the scenario's settings (e.g.\ fog density, sun altitude) or the run's
measurements (e.g.\ the initial separation); the assumption is recorded as
held, broken or not measured, and a failure is spurious only if an assumption
was broken. Clauses outside this form are kept as free text for the engineer.
Following Zave and Jackson~\cite{zave1997}, assumptions over the system under
test are rejected, so failures cannot be defined away by constraining the
system's own behaviour. Each assumption further receives a verdict
(load-bearing, not load-bearing, untested, insufficient data), answering
empirically whether it was needed.

\paragraph*{F3 Requirement-driven, reproducible test generation (stage 2)}
Besides grammatical evolution, \tool{} enumerates the scenario space
exhaustively, so failure rates are not biased by the search. Parameter ranges
deliberately exceed the stated assumptions, so that every assumption is
exercised. Scenarios carry global identifiers and seeds, giving paired
comparisons across system variants. Before cluster time is spent, \tool{}
compiles every scene against the map and verifies the cluster's code copy; it
then streams partial results as the job progresses.

\paragraph*{F4 A single oracle (stages 2, 3)} One path-aware safety rule (a
5\,m margin while the pedestrian is in the vehicle's path) scores runs in the
simulator's monitor, in the search's fitness, and in the analysis. A search
therefore cannot exploit an oracle the analysis does not use; in our runs, a
search did exploit the earlier, path-agnostic rule.

\paragraph*{F5 Automated failure analysis (stage 3)} \tool{} classifies every
run into a failure taxonomy (never detected, detected too late, brake released,
stopped too close, passed but stalled, among others) and separates test defects
from encounters. It checks REAL's catalogue obstacles against the evidence,
proposes obstacles of its own (e.g.\ \textsc{DetectionTooLate}), and surfaces
unnamed candidates from any factor that raises the failure rate.

\paragraph*{F6 Clustering of failure traces (stage 3)} Following the direction
set out in REAL~\cite{real2026}, \tool{} automates failure grouping by
clustering traces, adapting the approach of DYNASTO~\cite{dynasto2026}. \tool{}
encodes each run as the sequence of 13 events it exhibits (e.g.\ first
detection, braking released while moving, pedestrian enters the lane), one
\emph{set} of events per step, truncated at the closest approach so that runs
ending alike align regardless of duration. Runs are compared by an edit
distance whose substitution cost is the Jaccard distance between event sets,
normalised by length; the repeats of a scenario are collapsed to their medoid.
To make the choice of
hierarchical clustering parameters explicit, linkage and cut threshold are
swept and scored by adjusted Rand index against the rule-based
failure types, with a permutation null. For review, a behaviour map places
every encounter by classical multidimensional scaling and states its own
faithfulness (variance captured, rank correlation with the true distances);
clusters are cut at the highest silhouette, and cuts below 0.25 are labelled
neighbourhoods rather than failure types~\cite{kaufman1990}.

\paragraph*{F7 Human triage with automatic support (stage 4)} The review
interface walks the engineer through five steps, linking every number to its
simulations and replaying any run from its trace. For each obstacle the engineer
decides whether the world, the requirement or the system accounts for the
failures; the tool's suggestions remain unconfirmed until a person selects them,
car-side changes are refused as domain assumptions, and signing is blocked until
every item is decided.

\paragraph*{F8 Traceable refinement (stage 4)} Decisions are recorded and
validated as an artefact, and $R_1$ is written only if it parses, with every
change labelled [S], [D], [R], or [T] (a defect in the test itself). Domain and
requirement changes are re-diagnosed on the stored traces in seconds; system
changes start a new run that records its parent and the decision that caused
it. The hashes of the scene template and grammar, recorded in every run, flag
hidden changes between
iterations.

% F9 <- inventory F14-F15, F19-F20, F24-F27, F90-F91
\paragraph*{F9 Traceability between text, model and evidence (stages 1, 3)}
\tool{} presents each requirement as a single artefact with two synchronised
views. In the editor, every quoted phrase adopts the colour of its element in
the KAOS goal model (module, domain property, soft goal), and selecting a
phrase highlights the corresponding node. After execution, \tool{} judges each
part of the requirement individually (failed, broken, not checked) and marks
the verdict in both views. Obstacles are attached to the goals they obstruct;
selecting one highlights its chain from obstacle through obstructed goal to
responsible module. Every reported count links to the simulations behind it.

% F10 <- inventory F36, F92-F95
\paragraph*{F10 Explaining an individual failure (stage 3)} Before presenting
aggregates, \tool{} grounds the diagnosis in a concrete failure: the scenario
ranked worst for the system. A to-scale, top-down replay is reconstructed from
the recorded trace, overlaying the safety margin, in-path detections and brake
commands, and the decisive moments (first detection, first brake, closest
approach) are one action away. The engineer can switch between runs of the
scenario to contrast a failure with a pass, and inspect a camera recording of
a separate run of the same scenario.

% F11 <- inventory F83, F98
\paragraph*{F11 An adjustable behaviour map (stage 3)} \tool{} exposes the
clustering of F6 as an interactive map rather than a fixed partition. Each
encounter is a point encoding its outcome and cluster; the runs supporting the
selected obstacle are highlighted, and any point opens its trace. The engineer
moves between coarser and finer cuts, each annotated with its silhouette, and
thus decides how much structure to read into the data. Each cluster is
summarised by its outcomes, its over-represented settings and a representative
run.

% F12 <- inventory F43-F50, F126, F146-F147
\paragraph*{F12 Closing the loop from the interface (stages 2, 4)} Testing a
requirement takes a single action from the editor: \tool{} reuses the results
of an identical requirement if they exist, and otherwise validates the
requirement and the cluster, submits the job, and streams partial results as
they arrive. During revision, the effect of domain and requirement changes is
re-judged immediately on the stored runs (``$k$ of $N$ failures still count
against the system''). After signing, the refined requirement is executed as
the next iteration in one action, and its report is linked to its
predecessor.

% F13 <- inventory F12, F85-F89, F180
\paragraph*{F13 A self-explaining interface} A workflow tracker locates the
engineer in the loop, and per-page guided tours and contextual instructions
state what to inspect and what to decide. Tool output and human judgement are
visually distinct throughout. A built-in reference manual, REAL-DSL(5), in the
style of Unix manual pages, documents the syntax, the measurable quantities and
the available modules, with every example validated against the parser and the
diagnostics for common mistakes. The interface meets WCAG~2.2 AA contrast,
verified by automated tests, and never encodes meaning by colour alone.

% F14 <- inventory F150-F166
\paragraph*{F14 A built-in evaluation instrument} A study configuration of
\tool{} hides the tool's verdicts, assigns each participant an anonymous
identifier, and records every decision, rationale and interaction. Signed
sessions are uploaded to storage that participants can append to but not read.
This makes the human decision step itself an object of study: whether
engineers reach the same decisions from the same evidence.
